\documentclass[12pt, twoside]{article}
\input{../preamble}

\fancyhead[LE]{\thepage}
\fancyhead[RO]{Markscheme}
\fancyhead[LO]{La Scuola d'Italia / Dr.\ Huson / IB Math \\* 7 May 2026}

\begin{document}
\subsubsection*{7.7 Pre-Quiz: Logarithms --- Markscheme}

\begin{enumerate}[itemsep=0.3cm]

%--- Problem 1: AA SL 25M P1 TZ1 Q1 (qb_id 25M.1.SL.TZ1.1) — express in form ln k ---
\item[1.] [Maximum mark: 5]

\textbf{(a)} $\ln 3 + \ln 4 = \ln(3 \cdot 4) = \ln 12$. \hfill \textbf{A1}

\medskip

\textbf{(b)} $3\ln 2 = \ln 2^{3}$ \hfill \textbf{A1}

$= \ln 8$. \hfill \textbf{A1}

\medskip

\textbf{(c)} $-\ln\dfrac{1}{2} = \ln\!\left(\dfrac{1}{2}\right)^{\!-1}$ \hfill \textbf{A1}

$= \ln 2$. \hfill \textbf{A1}

\emph{Accept also $-\ln\tfrac{1}{2} = -(\ln 1 - \ln 2) = -(-\ln 2) = \ln 2$.}

\newpage

%--- Problem 2: AA SL 25M P1 TZ3 Q2 (qb_id 25M.1.SL.TZ3.2) — log_{10} expressions in p, q ---
\item[2.] [Maximum mark: 5]

$\log_{10} 2 = p$ and $\log_{10} 3 = q$.

\medskip

\textbf{(a)} $\log_{10} 24 = \log_{10}(8 \cdot 3) = \log_{10} 2^{3} + \log_{10} 3$ \hfill \textbf{M1}

$= 3\log_{10} 2 + \log_{10} 3$ \hfill \textbf{A1}

$= 3p + q$. \hfill \textbf{A1}

\medskip

\textbf{(b)} Change of base:
$\log_{3} 8 = \dfrac{\log_{10} 8}{\log_{10} 3} = \dfrac{3\log_{10} 2}{\log_{10} 3}$. \hfill \textbf{A1}

$= \dfrac{3p}{q}$. \hfill \textbf{A1}

\newpage

%--- Problem 3: AA SL 24M P1 TZ1 Q3 (qb_id 24M.1.SL.TZ1.3) — log_{10} a = 1/3 ---
\item[3.] [Maximum mark: 5]

$\log_{10} a = \dfrac{1}{3}$, where $a > 0$.

\medskip

\textbf{(a)} $\log_{10}\!\left(\dfrac{1}{a}\right)
= \log_{10} a^{-1} = -\log_{10} a$ \hfill \textbf{A1}

$= -\dfrac{1}{3}$. \hfill \textbf{A1}

\medskip

\textbf{(b)} Change of base:
$\log_{1000} a = \dfrac{\log_{10} a}{\log_{10} 1000} = \dfrac{\log_{10} a}{3}$. \hfill \textbf{A1}

$= \dfrac{1/3}{3}$ \hfill \textbf{A1}

$= \dfrac{1}{9}$. \hfill \textbf{A1}

\newpage

%--- Problem 4: AA SL Exemplar P1 Q2 (qb_id EXN.1.SL.TZ0.2) — solve 2 ln x = ln 9 + 4 ---
\item[4.] [Maximum mark: 5]

$2\ln x = \ln 9 + 4$, with $x = p\, e^{q}$, $p, q \in \mathbb{Z}^{+}$.

\medskip

$2\ln x - \ln 9 = 4$. \;Use $m\ln x = \ln x^{m}$: \;$\ln x^{2} - \ln 9 = 4$.
\hfill \textbf{M1}

Use $\ln a - \ln b = \ln\dfrac{a}{b}$:
$\ln\dfrac{x^{2}}{9} = 4$. \hfill \textbf{M1}

Exponentiate:
$\dfrac{x^{2}}{9} = e^{4} \;\Rightarrow\; x^{2} = 9e^{4}$. \hfill \textbf{A1}

Since $x > 0$:
$x = \sqrt{9e^{4}} = 3e^{2}$. \hfill \textbf{A1}

So $p = 3, \; q = 2$. \hfill \textbf{A1}

\emph{Accept any valid method, including expressing $4 = 4\ln e$ first
to combine all terms as a single log before exponentiating.}

\newpage

%--- Problem 5: AA SL 25N P1 TZ3 Q6 (qb_id 25N.1.SL.TZ3.6) — solve 3 log_8(10x) - log_4 x = 1 ---
\item[5.] [Maximum mark: 5]

$3\log_{8}(10x) - \log_{4} x = 1$, \;$x > 0$.

\medskip

Change of base to base $2$: \;$\log_{8} y = \dfrac{\log_{2} y}{3}$, \;
$\log_{4} y = \dfrac{\log_{2} y}{2}$. \hfill \textbf{M1}

Substitute:
$3 \cdot \dfrac{\log_{2}(10x)}{3} - \dfrac{\log_{2} x}{2} = 1$,
\;so $\log_{2}(10x) - \dfrac{1}{2}\log_{2} x = 1$. \hfill \textbf{A1}

Apply power rule and quotient rule:
$\log_{2}(10x) - \log_{2} x^{1/2} = \log_{2}\!\left(\dfrac{10x}{\sqrt{x}}\right)
= \log_{2}(10\sqrt{x}) = 1$. \hfill \textbf{M1}

\emph{Power rule and quotient rule may be applied in any order; this M1
rewards either rewriting toward a single logarithm with the same base.}

Equate arguments: $10\sqrt{x} = 2$, \;so\; $\sqrt{x} = \dfrac{1}{5}$.
\hfill \textbf{A1}

$x = \dfrac{1}{25}$. \hfill \textbf{A1}

\newpage

%--- Problem 6: AA SL 24N P1 TZ1 Q8 (qb_id 24N.1.SL.TZ1.8) — f(x)=log_2(8x), inverse, transformations ---
\item[6.] [Maximum mark: 14]

$f(x) = \log_{2}(8x)$, \;$x > 0$.

\medskip

\textbf{(a)(i)} $f(2) = \log_{2}(8 \cdot 2) = \log_{2} 16$. \hfill \textbf{M1}

$= 4$. \hfill \textbf{A1}

\medskip

\textbf{(a)(ii)} $f\!\left(\dfrac{1}{8}\right)
= \log_{2}\!\left(8 \cdot \tfrac{1}{8}\right) = \log_{2} 1 = 0$.
\hfill \textbf{A1}

\medskip

\textbf{(b)} Set $y = \log_{2}(8x)$ and swap $x$ and $y$:
$x = \log_{2}(8y)$. \hfill \textbf{M1}

Rewrite as exponential:
$2^{x} = 8y$. \hfill \textbf{M1}

$y = \dfrac{2^{x}}{8}$ \hfill \textbf{A1}

$f^{-1}(x) = \dfrac{2^{x}}{8} \;\;\left( = 2^{x-3} \right)$.
\hfill \textbf{A1}

\emph{Accept either form; $\dfrac{2^{x}}{8}$ and $2^{x-3}$ are
equivalent.}

\medskip

\textbf{(c)} $f^{-1}(0) = \dfrac{2^{0}}{8} = \dfrac{1}{8}$.
\hfill \textbf{A1}

\medskip

\textbf{(d)} $y = f(4x^{2}) = \log_{2}(8 \cdot 4x^{2}) = \log_{2}(32x^{2})$.

Use addition rule for logs:
$\log_{2}(32x^{2}) = \log_{2} 32 + \log_{2} x^{2}$. \hfill \textbf{M1}

$= 5 + \log_{2} x^{2}$. \hfill \textbf{A1}

Use exponent property: $\log_{2} x^{2} = 2\log_{2} x$. \hfill \textbf{M1}

So $f(4x^{2}) = 5 + 2\log_{2} x$. \hfill \textbf{A1}

Compared with $y = \log_{2} x$: vertical stretch (dilation) by scale
factor $2$, then vertical translation $5$ units upward. \hfill \textbf{A2}

\emph{Award A2 for both transformations stated with the correct order.
Award A1 if the two transformations are correct but the order is not
specified or is reversed.}

\end{enumerate}
\end{document}
